% TOP_PROBLEM: {"problemId": "problem.arveson-douglas-homogeneous-quotient-essential-normality", "canonicalTitle": "Arveson-Douglas Essential Normality Conjecture for Homogeneous Quotient Modules", "releaseRank": 463, "primaryDomain": "analysis", "primaryDomainLabel": "Analysis", "displayStatus": "open_with_solved_subcases", "releaseStatus": "open_with_solved_subcases"}
% !TeX program = lualatex
% ATTEMPT_STATUS: unresolved
\documentclass[11pt]{article}
\usepackage[margin=25mm]{geometry}
\usepackage{fontspec}
\setmainfont{Latin Modern Roman}
\usepackage{amsmath,amssymb,mathtools}
\usepackage{unicode-math}
\setmathfont{Latin Modern Math}
\usepackage{xurl,hyperref}
\hypersetup{colorlinks=true,urlcolor=blue}
\setcounter{secnumdepth}{0}
\setlength{\parindent}{0pt}
\setlength{\parskip}{5pt}
\setlength{\emergencystretch}{3em}
\begin{document}
\section*{\texorpdfstring{Arveson-Douglas Essential Normality Conjecture for Homogeneous Quotient Modules}{Arveson-Douglas Essential Normality Conjecture for Homogeneous Quotient Modules}}\label{463}
\subsection{Definitions and mathematical statement}
The Drury--Arveson space $H_d^2$ is the reproducing-kernel Hilbert space on the unit ball of ${\mathbb{C}}^d$ with kernel
\[
 K(z,w)=\frac1{1-\langle z,w\rangle}.
\]
For a proper homogeneous ideal $I\subset{\mathbb{C}}[z_1,\ldots,z_d]$, let $[I]$ be its Hilbert-space closure and $Q_I=H_d^2\ominus[I]$. Let $S_j=P_{Q_I}M_{z_j}|_{Q_I}$ be the compressed coordinate multiplier. For $p>0$, a compact operator $T$ lies in the Schatten class $\mathcal S^p$ when $\sum_ns_n(T)^p<\infty$, where $s_n(T)$ are its singular values.

The assertion is
\[
 [S_i^*,S_j]\in\mathcal S^p
 \quad\text{for all }i,j\text{ and every }p>\dim_{{\mathbb{C}}} Z(I).
\]
Here $Z(I)$ is the affine zero variety. All finite $d$ and all proper homogeneous ideals, including nonradical ones, are included; radical or smooth-variety cases alone do not give the full result.
\par\smallskip\textit{Formulation and concept references:} \href{https://users.math.cas.cz/~englis/85.pdf}{[S1]}.

\subsection{Short English statement}
Do homogeneous quotient modules of the Drury--Arveson space have commutators with the predicted Schatten-class decay, even for nonradical ideals?
\subsection{Research attempt: the compression defect and a nonradical family}
\textbf{Status.} The assertion for an arbitrary homogeneous ideal is not proved. A complete argument is given for coordinate power ideals, including arbitrarily thick nonradical examples. The calculation isolates a sufficient degreewise estimate for a general ideal and identifies exactly where replacing an ideal by its radical loses information.

\textbf{Source audit.} The Main Theorem on page 4 of [S1] treats the vanishing module of a homogeneous variety smooth away from the origin. Its introduction explicitly distinguishes the geometric, radical-ideal version from the full ideal formulation. The family below includes nonradical ideals, but belongs to the elementary monomial setting already discussed in that introduction; no novelty is claimed. This source check does not turn the geometric theorem into a theorem for every nonradical ideal or every singular projective variety.

\subsubsection{1. Grading and the uncompressed calculation}
Let $H_k$ denote the homogeneous polynomials of total degree $k$. Expanding the reproducing kernel shows that
\[
 e_\alpha=\sqrt{\frac{|\alpha|!}{\alpha!}}z^\alpha,
 \qquad \alpha\in\mathbb N^d,
\]
is an orthonormal basis. Consequently, with $k=|\alpha|$,
\[
 M_j e_\alpha=\sqrt{\frac{\alpha_j+1}{k+1}}e_{\alpha+e_j},
 \qquad
 M_i^*e_\alpha=\sqrt{\frac{\alpha_i}{k}}e_{\alpha-e_i}\quad(k\geq1),
                                                               \tag{1}
\]
where the second expression is zero when $\alpha_i=0$, and every backward shift kills the constant vector.

For $i\ne j$, direct composition gives
\[
 [M_i^*,M_j]e_\alpha
 =-\frac{\sqrt{\alpha_i(\alpha_j+1)}}{k(k+1)}
       e_{\alpha-e_i+e_j}.                                  \tag{2}
\]
For $i=j$ it gives
\[
 [M_i^*,M_i]e_\alpha=\frac{k-\alpha_i}{k(k+1)}e_\alpha.
                                                               \tag{3}
\]
The maps of indices in (2) are injective on their nonzero domain. Hence their operator norm is the largest absolute coefficient, rather than a sum of coefficients. Since $\sqrt{\alpha_i(\alpha_j+1)}\leq k+1$, (2)--(3) imply
\[
 \|[M_i^*,M_j]|_{H_k}\|\leq k^{-1}\qquad(k\geq1).           \tag{4}
\]
The degree-zero block is finite dimensional and will never affect a Schatten threshold.

For homogeneous $I$, put $Q_k=Q_I\cap H_k$ and $J_k=[I]\cap H_k$. Each finite-dimensional $I\cap H_k$ is closed; orthogonality of the grading then gives
\[
 Q_I=\bigoplus_{k\geq0}Q_k,\qquad H_k=Q_k\oplus J_k.
\]
Invariance of $[I]$ under every $M_j$ implies coinvariance of $Q_I$ under $M_j^*$. This last fact is essential: forward multiplication can leave $Q_I$, whereas backward multiplication cannot.

\subsubsection{2. An exact formula for what compression adds}
Let $P=P_{Q_I}$ and define the leakage maps
\[
 A_{j,k}=(1-P)M_j|_{Q_k}:Q_k\longrightarrow J_{k+1}.
\]
Using $S_j=PM_j|_{Q_I}$ and inserting $P+(1-P)$ yields
\[
 [S_i^*,S_j]|_{Q_k}
 =P[M_i^*,M_j]|_{Q_k}-A_{i,k}^*A_{j,k}.                    \tag{5}
\]
Indeed, $S_i^*S_j=PM_i^*PM_j|_Q$, while $S_jS_i^*=PM_jM_i^*|_Q$ because $M_i^*Q\subset Q$. The omitted term in the first product is precisely $PM_i^*(1-P)M_j|_Q=A_i^*A_j$. In particular its sign in (5) is negative.

This formula explains why the estimate for the ambient shift does not automatically settle the quotient problem. The ambient commutator is already $O(k^{-1})$, but multiplication followed by projection onto the ideal introduces a new term. Boundedness of coordinate multiplication gives only $\|A_{j,k}\|\leq1$, which has no decay in $k$.

Here is a useful sufficient criterion, proved without any operator-ideal interpolation. Suppose $r\geq1$ and
\[
 \dim Q_k\leq C(k+1)^{r-1},\qquad
 \|A_{j,k}\|\leq C_j(k+1)^{-1/2}.                          \tag{6}
\]
Then $[S_i^*,S_j]\in\mathcal S^p$ for all $p>r$. To see this, (4)--(6) bound each degree-$k$ block by $C'(k+1)^{-1}$. The blocks preserve degree and are orthogonal in both domain and range. Truncating to degrees at most $N$ is finite rank, and the norm of the remaining direct sum tends to zero. Thus the operator is compact. Each block has at most $\dim Q_k$ singular values, so for every $p>0$,
\[
 \sum_{\ell}s_\ell([S_i^*,S_j])^p
 \leq C''\sum_{k\geq0}(k+1)^{r-1-p}<\infty
 \quad\text{if }p>r.                                    \tag{7}
\]
This argument works also when $0<p<1$; it does not use a triangle inequality for a Schatten quasi-norm. All summation is over the singular values of orthogonal finite-dimensional blocks.

For a general homogeneous ideal, the expected rank exponent in (6) is the usual Hilbert-function exponent, $r=\dim_{\mathbb C}Z(I)$. The conditional criterion itself only assumes the displayed rank bound; the proof of the concrete family below obtains that bound by direct counting and does not require a general Hilbert-polynomial theorem. The unproved general step in this approach is the leakage estimate, not the summation in (7).

\subsubsection{3. Complete verification for coordinate thickenings}
Fix $0\leq r\leq d$ and positive integers $a_{r+1},\ldots,a_d$. Consider
\[
 I=(z_{r+1}^{a_{r+1}},\ldots,z_d^{a_d}),                   \tag{8}
\]
with the empty list interpreted as $I=(0)$ when $r=d$. Its zero variety is the coordinate $r$-plane. If any exponent is greater than one, (8) is generally nonradical: for example $z_j\notin I$ but $z_j^{a_j}\in I$ when $a_j>1$.

The standard monomials in the quotient are exactly those satisfying
\[
 0\leq\alpha_j<a_j\quad(j>r),                            \tag{9}
\]
with no restriction on the first $r$ coordinates. Orthogonality of the monomials proves that these $e_\alpha$ form an orthonormal basis of $Q_I$. In particular, there is no choice of a possibly ill-conditioned algebraic basis hidden in this identification.

If $j\leq r$, multiplying a standard monomial by $z_j$ preserves (9), so $A_{j,k}=0$. If $j>r$, leakage occurs precisely on the boundary layer $\alpha_j=a_j-1$. On that layer (1) gives
\[
 A_{j,k}e_\alpha=\sqrt{\frac{a_j}{k+1}}e_{\alpha+e_j}.
\]
These output monomials are orthogonal, and on the other layers the leakage is zero. Therefore
\[
 \|A_{j,k}\|\leq\sqrt{a_j}\,(k+1)^{-1/2}.                \tag{10}
\]
If the layer is empty in small degrees, the norm is simply zero; (10) remains valid.

For $r\geq1$, let $\beta=(\alpha_{r+1},\ldots,\alpha_d)$ and write $|\beta|$ for its sum. Directly counting the unrestricted coordinates gives
\[
 \dim Q_k=
 \sum_{\substack{0\leq\beta_j<a_j\\ |\beta|\leq k}}
 {k-|\beta|+r-1\choose r-1}
 \leq \left(\prod_{j>r}a_j\right){k+r-1\choose r-1}.
                                                               \tag{11}
\]
The stars-and-bars formula used here follows by inserting $r-1$ separators among $k-|\beta|$ identical objects. The last binomial is bounded by a constant times $(k+1)^{r-1}$. Equations (10)--(11) establish (6), and hence prove exactly the requested conclusion $p>\dim Z(I)=r$ for every ideal (8).

If $r=0$, all coordinates have bounded exponent in the quotient, so its dimension is $\prod_j a_j$. Every commutator is then finite rank and belongs to every $\mathcal S^p$, $p>0$, as required. This handles the dimension-zero endpoint separately instead of inserting $r=0$ into an inappropriate infinite-degree estimate.

For example, in four variables the nonradical ideal
\[
 I=(z_3^2,z_4^3)
\]
has variety dimension two and at most $6(k+1)$ standard monomials of degree $k$. The only leakage maps are in coordinates three and four, with squared norms at most $2/(k+1)$ and $3/(k+1)$. Formula (5) then gives summability for all $p>2$, independently of the ambient dimension four. This is an actual test of the refined dimension in the statement, not just the weaker threshold $p>d$.

\subsubsection{4. What the exponent does and does not say}
The threshold in this family is not merely an artifact of counting, at least when $r\geq2$. Restrict to the subfamily $I=(z_{r+1},\ldots,z_d)$, so the quotient is $H_r^2$. By (3), on degree $k$ the eigenvalues of $[S_1^*,S_1]$ are
\[
 \frac{k-\alpha_1}{k(k+1)}.
\]
For all standard monomials with $\alpha_1\leq k/2$, these eigenvalues are at least $1/[2(k+1)]$. There are at least $c_r k^{r-1}$ such monomials for $k$ sufficiently large. For $r=2$ this follows by choosing $\alpha_1$ from $0$ to $\lfloor k/2\rfloor$. For $r>2$, each such choice has
${k-\alpha_1+r-2\choose r-2}\geq c'_r k^{r-2}$ completions. Consequently
\[
 \sum_\ell s_\ell([S_1^*,S_1])^r
 \geq c\sum_{k\geq k_0} k^{r-1}(k+1)^{-r}=\infty.         \tag{12}
\]
Thus one cannot generally replace $p>r$ by $p\geq r$. When $r=1$, the unthickened quotient is the one-variable Hardy shift, whose self-commutator is rank one. That exceptional endpoint does not invalidate (12), which was explicitly proved only for $r\geq2$.

Replacing (8) by its radical retains the same variety but changes the Hilbert space. For instance $I=(z_2^2)\subset\mathbb C[z_1,z_2]$ has two infinite layers $\alpha_2=0,1$, whereas $\sqrt I=(z_2)$ has one. The difference is infinite dimensional, despite the equality of zero sets. It therefore cannot be dismissed as a finite-rank perturbation without an additional operator argument. Our proof controls the second layer through its shift weights; it does not remove that layer algebraically.

Likewise, a general linear change of polynomial variables need not be unitary on Drury--Arveson space. Coordinate orthogonality in (1) is preserved by unitary changes of variables, but not by arbitrary badly conditioned linear maps in a way that makes (10) automatic. The explicit calculation should not be extended to a new ideal simply by saying that generators can be put in normal form.

\subsubsection{5. Verification and the first remaining gap}
The accompanying finite script enumerates standard monomials for $(z_3^2,z_4^3)$, checks (11) against enumeration, and builds the degree blocks of (5) independently from the compressed shift formulas. It verifies the compression sign and the leakage norm on bounded degrees. Those calculations check indexing and algebra, not the infinite Schatten sum; (7), (10), and (11) prove that sum.

The approach now has a concrete stopping point. For arbitrary homogeneous $I$, show that the compression defect has sufficient decay and multiplicity control. A bound $\|A_{j,k}\|=O(k^{-1/2})$ together with the Hilbert-function rank estimate would suffice, but has not been established here. A weaker estimate might still suffice if the large singular values occupy sufficiently small subspaces; merely knowing the dimension of $Z(I)$ does not prove either assertion. The coordinate-power theorem, the endpoint obstruction, and formula (5) are complete partial results. The full conjecture remains unresolved by this attempt.

\subsection{Sources}
\begin{itemize}
\item[S1]Englis and Eschmeier, Geometric Arveson-Douglas conjecture.
\url{https://users.math.cas.cz/~englis/85.pdf}.
\textit{Formulation source.}
\end{itemize}
\end{document}
